% Section 2: Related Work
\section{Related Work}
\label{sec:related-work}

\paragraph{Cost-aware model routing.}
\textsc{FrugalGPT}~\citep{chen2024frugalgpt} uses cascade-based selection among API models to reduce cost while preserving quality.
\textsc{RouteLLM}~\citep{ong2025routellm} learns routers from human preference data for strong/weak model routing.
\textsc{HybridLLM}~\citep{ding2024hybridllm}, \textsc{RouterDC}~\citep{chen2024routerdc}, and \textsc{LLMRank}~\citep{agrawal2025llmrank} pursue similar cost--quality objectives with classifier, contrastive, or ranking-based gating.
Recent surveys and unified formulations further systematize routing and cascading trade-offs~\citep{moslem2026routing_survey,dekoninck2025cascade_routing}.
These systems select among \emph{model variants} or API tiers; we instead route among
heterogeneous \emph{agent strategies}---direct prompting, chain-of-thought, ReAct-style
tool use, and multi-agent coordination---that differ in procedural structure as well as
token cost.

\paragraph{Query difficulty and workflow adaptation.}
Difficulty-aware orchestration has moved from model selection toward workflow depth.
\textsc{DAAO}~\citep{su2026daao} estimates query difficulty with a VAE and adapts multi-agent workflow depth and LLM assignment per query.
\citet{tang2025task_complexity_mas} formalize when multi-agent systems help as task depth and width grow.
\citet{zhu2025llm_already_knows} estimate LLM-perceived difficulty from hidden states without generating outputs.
\textsc{STRIDE}~\citep{asthana2025stride} and \textsc{effGen}~\citep{srivastava2026effgen} provide complementary design-time modality selection and small-model agent frameworks.
Our QCE representation targets the same pre-inference decision point but uses planner-derived procedure structure rather than latent difficulty scalars or hidden-state value functions alone.

\paragraph{Multi-model fusion and ranking.}
LLM-Blender~\citep{jiang2023llmblender} ranks and fuses candidate generations from multiple
models, showing that the best pipeline often depends on the query and that multiple models
can remain viable for the same input.
Its design is largely \emph{post-hoc}: several models run (or their outputs are pooled),
then a ranker or fuser selects or combines responses.
We study \emph{pre-inference} routing that executes a single strategy per query while still
preserving complementarity among correct strategies during training via cost-soft targets
$p^\star$ (\S\ref{sec:method-supervision}).

\paragraph{Agentic reasoning strategies.}
Chain-of-thought prompting~\citep{wei2022cot}, ReAct~\citep{yao2023react}, and multi-agent frameworks such as AgentVerse~\citep{chen2024agentverse} define the strategy space we route over.
Rather than proposing a new single strategy, we study how to choose among existing strategies before inference using query complexity signals.

\paragraph{Graph-structured routers.}
GraphRouter~\citep{feng2025graphrouter} constructs a query--model graph and applies graph
neural message passing to predict which model to invoke.
AgentRouter~\citep{zhang2025agentrouter} and related entity-centric graphs extend this idea to multi-agent assignment.
Such methods encode \emph{affiliation} between queries and candidates but typically do not
model anticipated \emph{solution procedure}---subtask depth, tool dependencies, or
orchestration structure.
AAR derives routing features from an LLM planner's procedure DAG and a compact
complexity vector $\mathbf{z}(q)$ (\S\ref{sec:method-qce}--\ref{sec:method-representation}),
complementing semantic embeddings rather than replacing execution with graph propagation at
inference.

\paragraph{Positioning.}
RouteLLM, LLM-Blender, and GraphRouter are the closest published baselines for
query-conditioned LLM orchestration and inform our experimental design
(\S\ref{sec:experiments-setup}).
RouteLLM motivates our utility-based regret metric; LLM-Blender motivates soft supervision
over multiple successful candidates; GraphRouter motivates structured features beyond flat
text embeddings.
Our routing benchmark evaluates these ideas in a unified four-strategy setting with oracle
supervision and executed benchmark transfer (Tables~\ref{tab:hard-vs-soft}--\ref{tab:benchmark-transfer-em}).
Implementation notes and extended comparisons appear in
Appendix~\ref{app:related-extended}.
